\section{Study Design}

\subsection{Overall Design}

This is a prospective, open-label, single-arm, first-in-human study that combines
a Phase 1 standard 3+3 dose-finding evaluation of perioperative daraxonrasib
(RMC-6236) with an early feasibility evaluation of the on-premises large language
model (LLM) directed eight-arm robotic pancreaticoduodenectomy (Whipple) system,
a significant-risk Class II collaborative device. The study is conducted under a
combined IND/IDE submission: the drug component proceeds under the investigational
new drug pathway and the device component under the investigational device
exemption pathway, with a single integrated protocol, a single informed consent
that carries the Physical AI opt-out, and a single safety governance structure.
The central hypothesis is that the combined intervention can be delivered with an
acceptable early safety profile and can complete the assigned operative tasks
under continuous human oversight while a recommended Phase 2 dose of daraxonrasib
is identified, supporting progression to a later-phase efficacy study.

Enrollment is planned for up to $n=18$ participants assigned across three
daraxonrasib dose levels (DL1 through DL3) by the 3+3 rule with staggered
(sentinel) enrollment, so that no two participants in a cohort undergo surgery
within the same DLT-evaluation interval until the supervising surgeon and the
data safety monitoring board (DSMB) have reviewed the preceding case. A mandatory
Phase 0 simulation-validation gate precedes any patient-facing robotic activity
(\S4.3). The device operates within a phase-graduated staged-autonomy model: in
this Phase 1 protocol the system runs only at the clinician-controlled
(teleoperation or supervised) level and the supervised semiautonomous level, both
under continuous human oversight with an immediate manual override; full
autonomy is prohibited at this phase consistent with 21 CFR $\S$312.21(e), and
any future move to higher autonomy is reserved for a later phase and gated by
independent safety review, a Unified Safety Level (USL) reassessment, and DSMB
concurrence \cite{ecfr2026part812}. Figure~4 presents the staged-autonomy model.

\begin{mermaidfig}
\node[mmin]  (st1) at (-0.15,0)     {Stage 1\\Clinician-controlled\\teleoperation /\\supervised;\\full autonomy\\prohibited \S312.21(e)};
\node[mmdec] (r1)  at (5.2,0)   {Independent safety\\review + USL\\reassessment};
\node[mmstep](st2) at (10.75,0)   {Stage 2\\Supervised\\semiautonomous;\\Class II, continuous\\oversight, immediate\\manual override};
\node[mmdec] (r2)  at (10.75,-3.4){Accumulated safety\\data + DSMB\\concurrence};
\node[mmgoal](st3) at (5.2,-3.4){Stage 3\\Higher autonomy\\after review\\(reserved,\\later phase)};
\draw[mmedge] (st1) -- (r1);
\draw[mmedge] (r1) -- (st2);
\draw[mmedge] (st2) -- (r2);
\draw[mmedge] (r2) -- (st3);
\end{mermaidfig}
\vspace{-0.5cm}
\figcaption{Figure 4. Phase-graduated staged-autonomy model. Autonomy increases
only through gated review: this Phase 1 protocol operates at the
clinician-controlled and supervised semiautonomous levels (Class II, continuous
oversight with immediate manual override; full autonomy prohibited per 21 CFR
$\S$312.21(e)), and any move to higher autonomy requires an independent safety
review, a USL reassessment, and DSMB concurrence, for a later phase.}

\subsection{Scientific Rationale for Study Design}

A single-arm, open-label design is scientifically and ethically appropriate for
this first-in-human study, and a randomized comparator arm is not. The objectives
at this stage are to characterize early safety, to identify a recommended Phase 2
dose, and to demonstrate that a novel surgical system can perform its intended
function within hard cyber-physical limits; none of these is a comparative
efficacy question, and each is answered most efficiently and most safely in a
small, intensively monitored single-arm cohort. The early-feasibility pathway
recognized by the FDA for significant-risk devices is explicitly structured around
small first-in-human cohorts that generate the device-performance and safety
evidence needed before any controlled comparison is undertaken
\cite{fda2013earlyfeasibility}. Randomizing a vulnerable population to an
unproven combined drug-device intervention before that feasibility and safety
evidence exists would expose participants to risk without a commensurate
scientific gain and would be difficult to justify under the equipoise and
minimization-of-risk principles that govern human-subjects research.

The study population is a vulnerable one. Participants have a lethal malignancy
with a total five-year survival below 13\% \cite{Siegel2025}, are facing the most
technically demanding curative-intent abdominal operation, and may experience
therapeutic misconception about an LLM-directed robotic system. The single-arm
design concentrates oversight on each participant, permits sentinel staggered
enrollment with case-by-case DSMB review, and allows immediate pausing of
accrual on a prespecified safety signal, all of which protect participants more
effectively than a parallel randomized structure would at this stage. The 2025
Dutch nationwide cohort of 1000 robotic pancreaticoduodenectomies provides a
robust external human benchmark for the secondary surgical-quality endpoints
(\S3), so the absence of an internal control does not leave the surgical results
uninterpretable \cite{DutchCohort2025}. For these reasons the design is a
single-arm dose-finding and early-feasibility study, with a randomized controlled
comparison deferred to a later phase once safety, the RP2D, and device
feasibility are established.

\subsection{Justification for Dose}

\paragraph{Daraxonrasib dose and 3+3 escalation.}
Perioperative daraxonrasib is administered orally once daily. The escalation uses
exact doses anchored to the RASolute 302 pan-KRAS program, in which the
300~mg once-daily dose is the established recommended Phase 2 dose
\cite{rasolute302}. To preserve a conservative first-in-human margin in the
perioperative setting, the escalation begins below that reference dose and steps
up to it: DL1 is 160~mg once daily, DL2 is 220~mg once daily, and DL3 is the
300~mg once-daily reference dose. Three participants are enrolled at each level
and observed over a 28-day DLT window. If no DLT occurs among the first three,
the cohort escalates to the next level. If one of three experiences a DLT, the
cohort expands to six; a single DLT among six permits escalation, whereas two or
more DLTs at a level define that level as exceeding the MTD and fix the MTD at the
preceding level. The MTD is the highest level at which fewer than two of six
DLT-evaluable participants experience a DLT, and the RP2D is selected at or below
the MTD. Figure~5 shows the escalation scheme.

\begin{mermaidfig}
\node[mmin]  (dl1) at (0,0)      {DL1 160 mg QD\\starting dose,\\enroll 3};
\node[mmdec] (q1)  at (4.7,0)    {DLTs in first 3?\\(28-day window)};
\node[mmstep](dl1e)at (4.7,-2.6) {Expand DL1\\to 6};
\node[mmdec] (q1b) at (4.7,-5.2) {DLTs in 6?};
\node[mmstep](dl2) at (9.4,0)    {DL2 220 mg QD\\escalate, enroll 3};
\node[mmgoal](dl3) at (9.4,-2.6) {DL3 300 mg QD\\target dose,\\enroll 3-6};
\node[mmgoal](mtd) at (9.4,-7.2) {Declare MTD / RP2D\\at prior dose level};
\draw[mmedge] (dl1) -- (q1);
\draw[mmedge] (q1) -- node[above,font=\scriptsize, xshift=-0.1cm]{0 of 3} (dl2);
\draw[mmedge] (q1) -- node[right,font=\scriptsize, yshift=0.1cm]{1 of 3} (dl1e);
\draw[mmedge] (dl1e) -- (q1b);
\draw[mmedge] (q1b) -- node[below,font=\scriptsize, yshift=-0.025cm, xshift=-0.3cm]{1 of 6} (mtd);
\draw[mmedge] (q1b) -| node[below,font=\scriptsize,pos=0.25, xshift=-0.1cm]{$\geq 2$ of 6} (mtd);
\draw[mmedge] (dl2) -- (dl3);
\draw[mmedge] (dl3) -- node[right,font=\scriptsize, yshift=0.8cm]{tolerated} (mtd);
\end{mermaidfig}
\vspace{-0.45cm}
\figcaption{Figure 5.\ Daraxonrasib 3+3 dose-escalation scheme. Cohorts of three
advance when no dose-limiting\\toxicity is observed in the 28-day window, expand to
six on a single dose-limiting toxicity, and define\\the maximum tolerated dose and
recommended Phase 2 dose at the highest tolerated level, escalating\\from
DL1 (160~mg once daily) through DL2 (220~mg) to the DL3 reference dose (300~mg).}

\paragraph{Device ``dose'' analogue and the Phase 0 gate.}
The device has no pharmacologic dose; its dose analogue is the readiness evidence
required before any patient-facing activity, and it is escalated only by passing
the mandatory Phase 0 simulation-validation gate. Three conditions must be met
and documented before the first robotic case. First, the system must achieve a
Unified Safety Level (USL) rating of at least 7.0 on the surgical-readiness scale.
Second, the Phase 0 campaign must complete at least 1000 simulated procedures
across at least two independent simulation frameworks, so that the readiness
finding is not an artifact of any single simulator. Third, the sim-to-real
fidelity must fall within prespecified limits: a trajectory gap of less than 2~mm
and a tip-force gap of less than 0.5~N between the simulation prediction and the
operative record \cite{ecfr2026part812}. These limits operate together with the
intra-operative hard caps (a $\leq 3$~ms cross-arm emergency-stop budget, per-arm
tip-force caps of $\leq 3$~N, a cumulative cross-arm cap of $\leq 18$~N, and
vascular no-fly gating around the superior mesenteric and portal veins), so that
the device dose analogue is met only when both the simulation evidence and the
real-time safety envelope are demonstrably in place.

\subsection{End of Study Definition}

An individual participant is considered to have completed the study at the last
visit or procedure specified in the Schedule of Activities (\S1.3), which is the
24-month overall-survival assessment, or at the participant's death,
documented disease progression precluding further follow-up, or withdrawal of
consent, whichever occurs first. The end of the study as a whole is defined as the
date of the last scheduled assessment for the last participant, which is the last
participant's 24-month overall-survival assessment. After that date, no further
study-specific procedures are performed and the database is locked for the primary
analysis. Participants who discontinue the investigational intervention for a
protocol-defined reason (\S7) remain in long-term follow-up for survival to the
24-month endpoint unless they withdraw consent for follow-up.
